% ===== §9 The Political Economy of Beauty =====
\section{The Political Economy of Beauty: The Contested Remainder}
\label{sec:pe-beauty}

\noindent
The aesthetic perception of contingency and the spaciousness it requires have been presented, so far, as goods of the intimate relation, private capacities for a private flourishing. This section shows that they are nothing so contained. The capacity to perceive the beautiful, and the emptiness that perception requires, have a political-economic weight: they are forces in the world of capital, sought and fought over, and the contest over them is one of the quiet fronts of the struggle between the administered world and what resists it. The argument proceeds through a concrete case, the costly emptiness of a luxury department store, read first as the capture of the aesthetic by capital, then as the site of a remainder capital cannot seize, and finally as one point on a spectrum of such contests. The section closes by drawing from the case an emancipatory thesis: that the cultivation of aesthetic perception is a micro-politics of de-alienation, the building of a capacity that capital may use but cannot confiscate.

% ============================================================
\subsection{The costly emptiness of Ginza Six}
\label{sub:pe-gsix}

Consider a fact that ought, on a narrowly economic logic, to be impossible. Ginza Six is among the most expensive retail real estate in Tokyo, and Tokyo's is among the most expensive in the world; every square metre of its floor space commands a rent that makes its use a matter of the most careful calculation. And yet, on its sixth floor, in the heart of this most costly of spaces, Ginza Six gives over a large expanse to emptiness: a broad, open, uncrowded hall in which, periodically renewed, stands a major work of contemporary art, a vast installation, a hanging sculpture, a piece by an artist of the first rank, and the space is free, open to anyone, selling nothing. On the logic of maximal rent-extraction this is madness; the same floor area, partitioned into boutiques, would yield a fortune in rent that the empty hall forgoes. The store deliberately, expensively, leaves a great space empty and fills it with free beauty. The fact is not unique to Ginza Six, the costly atrium, the foundation gallery, the corporate plaza given to public art, are familiar features of the high-capital landscape, but Ginza Six states it with peculiar clarity, because the contrast between the price of the space and the gratuity of its use is so stark. Why does capital, which wastes nothing, pay so dearly for emptiness and beauty?

% ============================================================
\subsection{The first reading: capital's capture of the aesthetic}
\label{sub:pe-capture}

The first answer is that this is the most refined calculation rather than madness, and that the apparent gift is in truth the most sophisticated form of extraction. The empty hall and its free beauty are an instrument of it rather than a forgoing of profit, working by a logic more advanced than the rent it appears to waste. The beauty draws people in, it is a lure, an attraction that brings footfall the boutiques alone could not. It lengthens their stay, and lengthened stay is increased spending. It confers upon the whole establishment an aura of culture, refinement, and generosity that raises the brand and justifies the prices throughout. And in the age of the photograph, it manufactures spectacle: the free artwork is endlessly photographed and circulated, a vast unpaid advertisement performed by the visitors themselves, the empty hall converted into the most efficient marketing the money could buy. On this reading the emptiness is the most expensive advertising space in the building, and the free beauty is bait; the gift is a transaction in which the visitor, drawn by beauty and moved to share it, performs unpaid promotional labour and is induced to spend. This is the culture industry in its perfected form, the moment Adorno foresaw when the aesthetic, which ought to resist instrumental reason, is itself wholly instrumentalised, when even beauty, even the contemplative emptiness that seems the very opposite of commerce, is captured and made to serve the logic of the market \citep{adorno1997,debord1994}. The spectacle, in Debord's sense, has annexed the aesthetic; the beholding of beauty has become a moment in the circulation of capital. On this reading the paper's own prized capacities are revealed as already colonised: the aesthetic perception of contingency, the love of spacious beauty, are exactly the susceptibilities that capital has learned to farm.

% ============================================================
\subsection{The second reading: the irreducible remainder}
\label{sub:pe-remainder}

The first reading is true, and it is not the whole truth, and the dialectical task is to see what it leaves out without denying what it captures. For consider what actually happens to a particular person who, coming to the store for some errand of consumption, rounds a corner on the sixth floor and is met, unexpectedly, by the great open space and the work of art standing in it. Something happens in them. They stop; they are, for a moment, taken, the purposive errand suspended, the gaze caught, a small genuine experience of beauty undergone. And this experience is \emph{real}, whatever the motive of the space that occasioned it. The store's calculation does not reach into the experience and hollow it; the person is genuinely stopped, genuinely moved, genuinely, for that moment, released from the purposive consumption that brought them, into the disinterested beholding of something beautiful for its own sake. Here is the remainder the first reading cannot capture: the aesthetic experience, once it occurs, has a reality and a value that exceed and escape the instrumental frame that produced it. This remainder is, in the precise vocabulary of the philosophy of beauty developed earlier (\xref{sub:ap-antinomy}), the moment of the \emph{real} in the aesthetic experience, the kernel that resists symbolisation, and therefore resists the symbolic capture by which capital would convert the experience wholly into a commodity. Capital built the container for its own ends; but what happens in the container, in the person who is genuinely stopped by genuine beauty, is not exhausted by those ends, and cannot be, because the experience of beauty is a genuine suspension of the purposive, and a genuine suspension of the purposive is precisely what serves no purpose, what stands outside the means-end logic even when that logic arranged for it to occur.

This is the deep point, and it is the paper's whole argument returning in the register of political economy. The aesthetic experience is, by its nature, the beholding of something for its own sake rather than for its use; it is the moment in which the purposive is suspended and something is met as gratuitously there. Such a moment cannot be fully instrumentalised \emph{from within}, because the moment it is genuinely undergone it is, by its very structure, a stepping-outside of instrumentality, a small piece of non-alienated experience occurring in the very heart of the apparatus of alienation. Capital can build the container, can arrange the occasion, can profit from the footfall and the photographs; but it cannot make the person's actual experience of the beauty be itself an instrumental experience, for then it would not be an experience of beauty at all but merely another transaction, and the lure would fail. Capital needs the experience to be genuine in order to profit from it, a fake beauty would draw no one, and the genuineness it needs is exactly the remainder that escapes it. The store requires that the visitor really be stopped, really be moved; and in really being stopped and moved, the visitor has a moment of genuine non-instrumental experience that the store's instrumentality cannot reach into and convert. The remainder is structural, not accidental: capital's capture of the aesthetic depends on leaving intact precisely the non-instrumental core that defeats capture.

% ============================================================
\subsection{The dialectic: beauty as a contested field}
\label{sub:pe-dialectic}

The two readings do not cancel; they compose a dialectic, and the dialectic is the truth of the matter. Beauty, in the world of capital, is neither simply liberating nor simply captured; it is a \emph{contested field}, a site of struggle in which capture and remainder are locked together, each requiring and resisting the other. The first reading establishes that the aesthetic is real political-economic force, that beauty and spacious emptiness do genuine economic work, draw real footfall, command real rent, justify real premiums; capital's willingness to pay dearly for them is the hardest possible proof that they are powers in the world rather than private trivialities. The second reading establishes that this force, even as capital wields it, carries a remainder capital cannot confiscate, a non-instrumental core that, being the very thing that gives beauty its drawing power, cannot be converted into pure instrument without destroying the power it was meant to serve. Beauty is thus simultaneously the most capturable and the least confiscatable of goods: capturable because its force is real and capital can harness it, unconfiscatable because the force depends on a non-instrumental core that resists conversion. The contest between these is permanent and unresolvable, and it is fought, micro-cosmically, in every person who is genuinely stopped by beauty in a space designed to profit from their being stopped.

This connects the case directly to the dialectic of controlled contingency (\xref{sub:constr-controlled}). Ginza Six is a \emph{container built by capital}: its boundary is rigidly controlled for commercial ends, the space is a retail environment, arranged to sell, while its interior, the actual aesthetic experience it occasions, remains genuinely open, uncontrolled, the visitor's own. The structure is exactly that of the field: a controlled boundary enclosing an open interior. What the case adds is that the boundary's control and the interior's openness can serve \emph{different masters}, the boundary serving capital's commercial purpose, the interior serving the visitor's non-instrumental experience, and that this divergence is the very form of the contest. Capital controls the container; the remainder lives in the contents; and the political economy of beauty is the struggle over a field whose boundary one party holds and whose interior escapes them.

% ============================================================
\subsection{\texorpdfstring{\zh{无之以为用}}{Wu zhi yi wei yong}: emptiness confirmed by capital}
\label{sub:pe-emptiness}

There is a further and deeper point in the case, and it concerns emptiness specifically. What Ginza Six pays for, most strikingly, is \emph{emptiness} rather than only beauty, the open, uncrowded, spacious hall, the deliberate non-filling of the most fillable space. And in paying so dearly for emptiness, capital confirms, against its own filling logic, the ancient thesis of the eleventh chapter: that the use of the vessel is in its emptiness, that what is not there gives use \citep{laozi1963}. The most calculating filler of space, the logic that would partition every square metre into rentable use, is driven, by the discovery that emptiness draws, that spaciousness has a power density of filling lacks, to leave a great space empty, and thereby to confirm \zh{无之以为用} in the one register that cannot be accused of mysticism: the register of rent. Capital has discovered, and paid to confirm, that emptiness is of use. The Daoist metaphysics of the void, which the paper has carried throughout as the deepest ground of the field, receives here its most unlikely and most decisive proof: that even the logic of maximal filling is forced, by the nature of value, to acknowledge that the empty is of use, and to pay the highest price for it.

But the same case shows that emptiness, too, is contested. The spaciousness capital pays for is spaciousness instrumentalised, emptiness deployed as a lure, the void made to work for the filling it appears to refuse. The atmospheric emptiness of the luxury hall is not the same as the spaciousness of the field, even though it borrows its form; it is emptiness in the service of consumption, the void as ambience, designed to dispose the visitor to spend. And yet, the remainder again, the emptiness, even instrumentalised, still opens a real space in which a real suspension of the purposive can occur; capital's instrumental emptiness still, despite itself, affords moments of the genuine spaciousness it was built to exploit. Emptiness, like beauty, is a contested field: paid for by capital, deployed as ambience, and yet harbouring a remainder that escapes the deployment, a genuine openness within the instrumentalised one. The struggle over beauty is, at its root, a struggle over emptiness, over whether the space that is opened will be a space for genuine non-instrumental experience or merely an ambience engineered for consumption, and the two are perpetually entangled, the genuine void haunting the engineered one it is made to wear.

% ============================================================
\subsection{The spectrum of cases}
\label{sub:pe-spectrum}

Ginza Six is one point on a spectrum, and placing it among others clarifies what is and is not specific to it. Consider three cases, ordered by the degree to which the aesthetic is captured by capital.

At one end stands the \emph{unownable} aesthetic: the beauty of the mountain and the water, the landscape that the wayfaring tradition perceived (\xref{sub:ap-eastern}), the contingent beauty of the wildflower at the roadside. This beauty is, in its nature, resistant to capture, because it is not produced by capital, not enclosed, not sold; the mountain's beauty is free in a sense the gallery's is not, available to anyone with the trained perception to catch it, owned by no one and rentable by no one. Capital can sell access to the mountain, the ticket, the lodge, the guided tour, but it cannot own the beauty itself, which remains available to the unmonied wayfarer with the seeing eye, and which the poorest perceiver with the richest perception possesses more fully than the wealthiest tourist who photographs it without seeing. This is the aesthetic at its most unalienated: a beauty no one made and no one can confiscate, requiring for its enjoyment nothing but the cultivated perception this paper has described.

In the middle stands the \emph{contested} aesthetic of Ginza Six: beauty produced and enclosed by capital, deployed for profit, and yet harbouring the irreducible remainder. Here the capture is real and the resistance is real, and the case is dialectical through and through, the gallery is neither the free mountain nor the pure trap, but the contested field in which both forces operate at once.

At the other end stands the \emph{extracted} aesthetic, the case in which capture is most nearly complete: the photogenic site engineered for the image, the ``Instagrammable'' location whose entire reality is its capacity to generate shareable photographs, the destination that exists to be photographed and has no other being. Here the aesthetic has been most thoroughly instrumentalised: the place is beautiful to be \emph{captured} rather than beautiful to be beheld, designed backward from the photograph, its every feature calculated for the image it will yield. And here the remainder is most nearly extinguished, not because the structure differs from the gallery's, but because the perception the place cultivates is the perception that does not behold but only captures, the gaze that does not undergo beauty but only records it for circulation. The extracted aesthetic is the triumph of dead material over living perception (\xref{sub:reproduction-deadlabour}): a place built for the photograph, undergone by no one, the beauty wholly converted into the image and the experience wholly evacuated. Yet even here the remainder is not quite zero, for a visitor with the cultivated perception might still, against the grain of the place's design, actually behold rather than capture, and have a genuine moment the engineered site did not intend. The spectrum has no endpoint of total capture, because the remainder, however reduced, is structural; but the extracted case shows how nearly it can be approached, and how the difference between the captured and the free turns, in the end, not on the place but on the perception brought to it.

% ============================================================
% ============================================================
\subsection{How beauty creates a value cycle: the mechanism}
\label{sub:pe-cycle}

The analysis so far has established that beauty is a real political-economic force, that it draws footfall, commands rent, justifies premiums, but it has treated this force statically, as a power beauty possesses, without showing the \emph{mechanism} by which beauty exercises it: how, dynamically, beauty creates and sustains a cycle of value. This mechanism must now be set out, for it is the link between the paper's phenomenology of aesthetic perception and its political economy, and it shows the political economy of beauty to be not a static fact but a dynamic cycle with the same structure as the relational reproduction the paper analysed at its centre.

The cycle has a definite point of ignition, and the point is the act of perception. Beauty is not a passive property awaiting valuation; it is an active power that \emph{seizes perception}, stops the passer-by, catches the gaze, draws the subject out of the purposive errand that brought them into the non-purposive beholding of something for its own sake (\xref{sub:ap-moments}, \xref{sec:aesthetic-perception}). This seizure is the ignition of the value cycle. Before it, the subject is a purposive agent moving through means toward ends; in it, the subject is arrested, captured, released for a moment into the aesthetic experience, and it is exactly this arrest, the conversion of a purposive passage into a non-purposive beholding, that beauty performs and that nothing else performs so well. The person stopped on the sixth floor of Ginza Six by the great space and the work standing in it (\xref{sub:pe-emptiness}) is the value cycle igniting: a purposive consumer seized, against the grain of their errand, into a genuine aesthetic experience. Beauty makes the subject perceive; the making-perceive is the first moment of the cycle.

The second moment is the generation of value in the perception itself, and here the three-register analysis specifies what the value is (\xref{sub:ap-threeregisters}). The seized perception is a genuine experience, the Real surplus of the beautiful thing met, invested with Imaginary captivation, raised toward the Symbolic of the shareable, and this experience is itself a value: not yet an economic magnitude, but a lived good, the non-instrumental experience that the subject undergoes and that has worth in the undergoing. This is the value beauty generates at the source: the aesthetic experience as a lived good. And it is, crucially, a value with the double character the dialectic established (\xref{sub:pe-dialectic}): it can be captured economically, converted into footfall, dwell-time, the premium the subject will pay to be where such experiences occur, and yet it carries the irreducible remainder, the non-instrumental core that the economic capture depends upon and cannot itself produce. The value beauty generates is at once capturable (it can be turned to economic account) and unconfiscatable (its economic capture requires the genuineness that escapes capture).

The third moment is the one that makes this a \emph{cycle} rather than a single event: the reproduction of the value, by which the aesthetic experience generates further perception and so renews the source. The aesthetic experience, once undergone, is reproduced, recounted, returned to, shared, made into a reason to come back and a reason others come (\xref{sec:reproduction}). And this reproduction feeds the source: the recounted experience draws new perceivers, the return visit renews the experience, the shared beauty propagates the seizure of perception to others, so that the value generated at the source flows back to generate more perception, more experience, more value. This is the value cycle proper: beauty seizes perception, perception generates value, value is reproduced, and the reproduction renews and amplifies the seizure of perception, a self-sustaining circulation in which beauty, by being perceived and the perception reproduced, generates a flow of value that returns to its source and augments it. The political economy of beauty is this cycle: not a static force but a dynamic circulation, ignited by aesthetic perception and sustained by its reproduction.

% ============================================================
\subsection{The good and the bad value cycle}
\label{sub:pe-goodbadcycle}

That beauty creates a value cycle does not yet say whether the cycle is good or bad, and the distinction, which the paper has drawn for every other cycle it has analysed (\xref{sub:reproduction-cycle}, \xref{sub:ev-dynamics}), applies here with full force and full discriminating power. The value cycle that beauty ignites can run as a good cycle, accumulating value and preserving the remainder, or as a bad cycle, extracting value and consuming the remainder until beauty itself is hollowed out.

The good value cycle is the one in which the reproduction of the aesthetic experience is \emph{nourishing}, in which each turn of the cycle deepens the experience, enriches the perception, and conserves the non-instrumental remainder that is beauty's living core. Here beauty is perceived, the perception generates genuine experience, the experience is reproduced in ways that renew rather than exhaust it, and the cycle accumulates value while keeping beauty alive: the place returned to that grows richer with each visit, the beauty whose reproduction deepens it, the cultivated perception that finds more in it each time. The value accumulates, in the deepening of experience, in the reputation that draws others, in the relational wealth of those who share it, and the remainder is preserved, because the reproduction is the kind that nourishes rather than extracts. This is the good cycle in the register of the political economy of beauty: a circulation that returns more than it takes, the holonomy of value accumulating along the path of nourishing reproduction (\xref{sub:reproduction-cycle}), beauty growing rather than diminishing as it is perceived and reproduced.

The bad value cycle is the one in which the reproduction is \emph{extractive}, in which each turn consumes the experience, depletes the perception, and hollows out the remainder until beauty is reduced to a spent commodity. Here beauty is perceived, but the perception is reproduced in ways that exhaust it: the experience consumed for its extractable value, the beauty photographed without being beheld, the place worn smooth by an extraction that takes the value and gives nothing back, until the remainder, the genuine non-instrumental core that drew perception in the first place, is consumed away, and beauty, hollowed, loses the very power by which it ignited the cycle. This is the Instagram extreme of the spectrum (\xref{sub:pe-spectrum}): the beauty extracted so thoroughly that nothing is left to behold, the experience consumed until the place that drew millions for its beauty no longer holds any, the remainder spent, the cycle collapsing as the beauty it fed upon is exhausted. The bad value cycle is self-undermining, for it consumes its own source: extracting the remainder that was the condition of beauty's drawing-power, it kills the goose, and the cycle that hollowed out its beauty finds, at last, that it has nothing left to circulate. The decisive variable, here as everywhere in the paper, is the manner of its reproduction rather than beauty itself, nourishing or extractive, conserving the remainder or consuming it, and the same beauty that, well reproduced, founds a good cycle of accumulating value is, badly reproduced, fed into a bad cycle that consumes it to nothing.

% ============================================================
\subsection{The value cycle as relational production}
\label{sub:pe-cyclerelational}

The deepest truth about the value cycle of beauty is that it is, at its root, a cycle of \emph{relational production}, and this returns the political economy of beauty to the paper's unified core (\xref{sec:creation}). For the value that beauty generates is not, fundamentally, a value residing in objects or accruing to capital; it is the value generated when a subject is drawn, through beauty, into a relation, with the beautiful thing, with the place, with those who behold it together, with the world the beauty opens, and undergoes, in that relation, the generation of a value that belongs to the relation. The aesthetic experience that ignites the cycle is the coming-to-presence of a relation between the perceiver and what is perceived; the value it generates is relational value; and the reproduction that sustains the cycle is the reproduction of that relation, the relational production analysed at the paper's centre (\xref{sec:relational-production}). Beauty creates a value cycle because beauty draws subjects into relations, and relations, reproduced, generate and circulate value: the value cycle of beauty is relational production seen in the register of political economy, the same bringing-forth of value by a relational subject that the paper has tracked under the names of creation, reproduction, and the good cycle, now seen as the engine that beauty ignites and sustains.

This is why the good value cycle and the bad correspond so exactly to the good and bad cycles of relational production (\xref{sub:rp-criteria}). The good value cycle nourishes the relation that beauty opens, preserves its generativity, conserves the remainder that is the relation's living core, it satisfies, in the economic register, the criteria of good relational production. The bad value cycle extracts from the relation beauty opens, consumes its generativity, hollows out the remainder, it is the bad relational production, the extractive cycle, in the economic register. And the remainder that capital cannot confiscate (\xref{sub:pe-remainder}) is, seen now, precisely the relational core of the aesthetic experience: the non-instrumental relation between perceiver and perceived that is the genuine source of beauty's value, which capital can build a container around and profit from but cannot itself produce, because it is generated only in the genuine relation that beauty opens and that no instrumental arrangement can manufacture. The political economy of beauty is, in its depths, the political economy of relational production: beauty creates value by drawing subjects into relations that generate value, and the cycle it ignites is good or bad according to whether that relational production is nourished or extracted, the same distinction, the same criteria, the same keystone, now seen to govern the circulation of value that beauty sets in motion.

% ============================================================
\subsection{The emancipatory upshot: aesthetic education as micro-politics}
\label{sub:pe-emancipation}

The spectrum yields the section's emancipatory thesis, and with it the political stakes of the whole aesthetic argument. What distinguishes the three cases, in the end, is less the place than the \emph{perception} brought to it, for the cultivated perception can find the unalienated remainder even in the contested gallery and the extracted site, while the uncultivated perception misses the free beauty even of the mountain, photographing the landscape without ever beholding it. The decisive variable is the perception, and the perception is cultivable. This is where the paper's aesthetics becomes a politics. To cultivate the aesthetic perception of contingency, the capacity to behold rather than capture, to be genuinely stopped by genuine beauty, to undergo the non-instrumental experience for its own sake, is to cultivate a capacity that capital can use but cannot confiscate: a capacity for non-alienated experience that, once possessed, can be exercised even within the apparatus of alienation, finding the remainder that escapes capture wherever beauty occurs, in the gallery built to profit from it and the site built to extract it and the mountain that no one built at all.

This is a \emph{micro-politics of de-alienation}. Alienation, in the classical analysis, is the condition in which one's activity and experience are are appropriated rather than one's own, instrumentalised, turned into means for ends not one's own \citep{marx1992}; and the administered world extends this condition across ever more of life, until even leisure, even beauty, even the contemplative suspension of purpose is captured and made to work. Against this, the cultivation of aesthetic perception builds a capacity for genuine non-instrumental experience that capital cannot fully appropriate, because its value lies precisely in the non-instrumentality that capture would destroy. Each genuine beholding of beauty, each real suspension of the purposive, is a small piece of de-alienated experience reclaimed, a moment in which one's experience is genuinely one's own, undergone for its own sake, escaping the means-end logic even when occurring in its midst. And this reclamation is available to anyone, requiring no wealth, no leisure, no apparatus, only the cultivated perception that this paper has argued can be developed, and developed, between two people, mutually, as the relation's own aesthetic education (\xref{sub:ae-cultivation}). The politics of perception that Rancière named (\xref{sub:ae-modern}) is here given its intimate and practical form: two people cultivating in each other the capacity to perceive the unalienated remainder, and so building between them a small commons of de-alienated experience that the administered world cannot enclose. The cultivation of beauty's perception is not a retreat from politics into private aesthetics; it is a micro-politics, the building of a capacity for freedom at the scale of a perception, exercised in the catching of a wildflower and the beholding of a remainder, available to all and confiscatable by none. \el{χαλεπὰ τὰ καλά}: the beautiful is difficult, and its difficulty is now seen to be also its political promise, for what cannot be reduced to a rule cannot be wholly captured by a system, and the un-programmable perception of the beautiful is, for that very reason, the seed of a freedom the administered world cannot finally administer.
